\subsection{The arbiter and the witness}
\label{sec:ha:enabled:arbiter}

Two instances of a group may act, and only one of them may act at a time. This
subsection says how the venue decides which one.

An instance cannot make this decision alone. When an instance loses contact with
its peer, it cannot distinguish a peer that has stopped from a peer it can no
longer reach. Both readings lead to a wrong action. If it promotes itself and the
peer is alive, both instances act and their state diverges. If it declines and
the peer is dead, neither instance acts and the venue stops trading. The second
failure is the safer one, and it takes longer to notice, because no process
crashes and nothing logs an error.

\textbf{A majority of three voters decides.} Each group has three voters: its two
instances and the arbiter pool. An instance acts only while a majority of the
three has granted it a lease that has not run out. Its own vote is one of the
three, so it needs a grant from one other voter: its peer, or the arbiter pool. A
voter grants a lease to one instance at a time, and promises not to grant another
until the lease period has passed. Two instances would each need two of the three
votes at the same moment, so at least one voter would have to have granted both,
and no voter does. Every message also carries the generation number of the
entitlement it was sent under, and a receiver refuses a message whose generation
has been superseded.

\textbf{Two arbiters run, and one witness runs alongside them.} The arbiter pool
votes through whichever arbiter is active. The two arbiters decide which of them
is active by the same rule, with the witness as their third voter: an arbiter is
active only while a majority of the two arbiters and the witness has granted it a
lease. Losing one arbiter moves no entitlement and interrupts nothing, which is
\reqref{R-0063}. The witness has no pair and no standby.
\Cref{sec:ha:enabled:arbiter:witness} states what the venue does about that.

\subsubsection{What is done about the witness being single}
\label{sec:ha:enabled:arbiter:witness}

\textbf{The witness holds no state, so losing it loses nothing.} It has no record
to recover and no history to lose. After it restarts it grants nothing for one
lease period, which is the whole of its recovery.

\textbf{Losing the witness alone changes nothing.} Two arbiters that can see each
other grant each other leases without it. The absence of the witness is felt only
when an arbiter is also lost: the surviving arbiter then holds only its own vote,
and it is active only while its existing lease lasts. The pool then has no active
arbiter until another voter returns. Component leaders are unaffected while their
peers are running, because each renews its lease with its peer.

\textbf{The witness is off every ordinary path.} Only the arbiters talk to it. No
order, report or component touches the witness.

\textbf{The witness runs where the arbiters do not.} A witness that shares a
failure domain with one arbiter gives little protection, because a single event
then removes that arbiter and the witness together, and the surviving arbiter
holds only its own vote. The witness therefore runs on a third machine, on a
different power supply, behind a different network switch, and ideally in a
different rack. Its usefulness comes from where it is deployed rather than from
what it does.

Three properties of the witness process make that placement cheap:

\begin{itemize}[nosep]
  \item \textbf{It needs no storage.} It holds nothing on disk, so it depends on
        no array and no filesystem.
  \item \textbf{It is off the hot path.} It needs no pinned cores, no low-latency
        hardware and no proximity to any other component, so it can run on the
        least capable machine in the deployment.
  \item \textbf{Its network route matters more than its speed.} The failure that
        hurts the venue is an arbiter being unable to reach the witness. The
        route to the witness should therefore fail independently of the route
        between the two arbiters.
\end{itemize}

The venue goes on trading when both arbiters and the witness are gone, for as long
as each group's leader can renew its lease with its peer. What it cannot do then is
replace a leader that fails: \cref{sec:ha:enabled:degraded} says what happens.

\subsubsection{What state it holds}
\label{sec:ha:enabled:arbiter:state}

The arbiter holds whether it is active and, while it is, one vote per group: the
instance it has promised its vote to, if any, and until when. It also holds the
highest generation granted in each group. It holds nothing about orders, sessions
or sequence numbers. That limit is deliberate: a component that decides which
instance may act must not itself hold state that can be wrong.

The witness holds its vote on which arbiter is active, and nothing else.

\subsubsection{Where the copy that outlives it lives}
\label{sec:ha:enabled:arbiter:durable}

The arbiter keeps nothing across its own restart. It does not need to know what it
promised before it stopped: it grants nothing for one lease period after starting,
by which time every such promise has run out. The same applies to an arbiter that
becomes active in place of the other, because it does not know what the other
promised.

The highest generation granted in each group is the one thing worth keeping, and
the active arbiter tells the passive one each time it rises. An arbiter that
becomes active therefore knows it, unless both arbiters have restarted.

\subsubsection{How a restarted instance becomes current}
\label{sec:ha:enabled:arbiter:recovery}

A restarted arbiter waits one lease period and then votes. Until then it refuses
every request, which is \reqref{R-0059}.

\subsubsection{What the member is told}
\label{sec:ha:enabled:arbiter:member}

Nothing. The arbiter sits off the path a member's orders travel, and a member
cannot tell that one exists. An arbiter's failures reach a member only through
the components that consult it, so the requirements here are stated in terms of
those components rather than of members.

\begin{requirement}{R-0058}{Two instances of a group are never both entitled to act}
  At any moment, at most one instance of a group shall hold the entitlement to
  act.
  \rationale{This is the property the arbiter exists to provide, and the rest of
  this subsection follows from it. Two instances acting on the same state
  diverge, and nothing afterwards can establish which of them was right.}
  \uncovered{no scenario asserts the property directly, only particular routes to it}
\end{requirement}

\begin{requirement}{R-0146}{An instance acts only while a majority of its three voters grants it a lease}
  An instance of a group shall act only while it holds a lease, not yet run out,
  from its peer or from the arbiter pool. It shall count each lease from the
  moment it asked for it, less an allowance for clocks that run at slightly
  different rates. A voter shall grant a lease to one instance at a time, shall
  grant nothing for one lease period after it starts unless it has kept a record,
  made during the same boot of its machine, of what it promised, and shall grant
  no generation below one it has already granted. An instance that has granted its
  peer a lease shall not ask to lead until that lease has run out.
  \rationale{An instance that loses contact with its peer cannot tell a dead peer
  from an unreachable one, so no rule it applies alone can be right in both
  cases. A majority of three can be: two instances would each need two of the
  three votes at the same moment, which would need a voter that has granted both.

  Each clause closes a way in which two instances come to act at once, and model
  checking showed each one is needed: without it, two instances act as leader at
  once within a few steps. Counting the lease from the request rather than the
  grant means the holder always believes its lease ends first. Waiting after a
  start covers a voter that has forgotten what it promised; a voter that kept a
  record has not forgotten, and so can resume at once, which is what lets a
  supervised restart keep the lead. The epoch rule stops
  a voter granting an older generation. An instance that granted its peer a lease
  must not then count its own vote for itself.}
  \covers{ha_test:8, ha_test:15, ha_test:39, ha_test:49, tla:lease-1-holder-counts-from-arrival, tla:lease-2-candidate-ignores-own-grant,
  tla:lease-3-restart-does-not-wait}
\end{requirement}

\begin{requirement}{R-0059}{An arbiter that does not know what was promised grants nothing until every promise has run out}
  An arbiter that has just started, or has just become the active arbiter, shall
  grant no lease to any component for one lease period.
  \rationale{An arbiter keeps nothing across its own restart, and an arbiter that
  becomes active does not know what the previously active one promised. During
  that period it cannot distinguish two situations: a group in which no instance
  holds a lease, and a group whose leader holds one it cannot see. Granting a
  lease is correct in the first situation. In the second it gives a group two
  instances entitled to act, which is the outcome the arbiter exists to prevent.

  Waiting one lease period settles it without any knowledge: by then every
  promise made before is certain to have run out. The leader of each group
  renews with its peer meanwhile, so the wait costs nothing while a pair has both
  instances running.}
  \uncovered{no scenario asks a freshly started arbiter about a group it has not been told of}
\end{requirement}

\begin{requirement}{R-0060}{Silence from an arbiter is not read as the absence of one}
  A component that asks and hears nothing shall ask again before concluding that
  no arbiter exists.
  \rationale{An arbiter that stays silent may be the passive one, or one waiting
  out a lease period under \reqref{R-0059}. A component that treats the silence as
  an absent arbiter and acts anyway creates the second leader the silence was
  preventing. A component therefore leads only on a grant, never on the absence
  of an answer.}
  \uncovered{no scenario has a component hear nothing from any voter and checks that it does not lead}
\end{requirement}

\begin{requirement}{R-0061}{Entitlement is held under a lease and lapses unless renewed}
  An instance's entitlement to act shall expire unless it is renewed, and an
  instance whose entitlement has expired shall stop acting at once.
  \rationale{A record that never expires states only that something was true at
  the moment it was written. Presence alone cannot be the test either: an
  instance that holds its connection open and stops renewing looks, from the
  outside, the same as a healthy one. An instance that keeps acting after its
  lease has run out is a second leader the moment a voter grants its peer a
  lease, so stopping is what makes the lease mean anything.}
  \uncovered{no scenario stops a leader renewing while it stays connected}
\end{requirement}

\begin{requirement}{R-0099}{Adjusting a machine's clock does not move an entitlement}
  The expiry of a lease shall be measured by a clock that only moves forward at a
  steady rate, and shall not be affected by the machine's notion of the time of
  day being corrected, stepped or set.
  \rationale{Time-of-day clocks are adjusted routinely, by a time service
  correcting drift, by a leap second, or by an operator. A lease measured against
  such a clock can expire while it is being renewed properly, which moves an
  entitlement for a reason unconnected to any component's health. An adjustment
  in the other direction is worse: a lease that should have expired goes on being
  honoured after its holder has gone.}
  \uncovered{no scenario adjusts a clock}
\end{requirement}

\begin{requirement}{R-0109}{An instance that is not making progress does not keep its entitlement}
  Holding an entitlement shall depend on the instance continuing to do the work
  the entitlement is for, and not only on its continuing to ask to keep it.
  \rationale{Renewing an entitlement requires no more than a running thread. An
  instance starved of processor time, or stalled in a way that leaves its timers
  running, renews on schedule and does nothing else. Neither presence nor renewal
  reports whether the instance is doing any work, so both report such an instance
  as healthy. The venue then stays with an instance that has stopped, and
  everything waits while nothing appears to be wrong.}
  \uncovered{no scenario starves an instance that keeps renewing}
\end{requirement}

\begin{requirement}{R-0115}{How long an entitlement lasts without renewal is a stated figure}
  The period after which an unrenewed entitlement lapses shall be a figure the
  venue states, and it shall be long enough that an instance doing its work is
  never mistaken for one that has stopped.
  \rationale{A period that is too short takes the entitlement away from a busy
  instance for being busy, which moves the venue's work to the other instance at
  the moment there is most of it. A period that is too long leaves the
  entitlement with an instance that has genuinely stopped, and nothing acts until
  it lapses. The figure alone shows neither error, so it has to be derived from
  how long the work takes.}
  \uncovered{no scenario tests the period, and none is stated to test}
\end{requirement}

\begin{requirement}{R-0062}{An instance already acting keeps its entitlement when another returns}
  Where an instance holds the entitlement and is still present, an instance
  rejoining the group shall be told to follow, whatever its configured name.
  \rationale{Preferring the configured primary at this point moves the
  entitlement for no reason but a name. It also moves the entitlement to the
  instance that has just failed, and away from one that is serving. The
  configured preference is a cold-start tie-break.}
  \covers{ha_test:24, ha_test:25}
\end{requirement}

\begin{requirement}{R-0063}{Losing one arbiter, or the witness, changes nothing}
  The failure of a single arbiter, or of the witness, shall not move any
  group's entitlement and shall not interrupt the venue.
  \rationale{Redundancy is what makes the arbiter available, so nothing needs to
  watch it. A design in which losing one arbiter mattered would need something to
  supervise the arbiter, and that supervisor would need supervising in turn.}
  \covers{ha_test:2, ha_test:4, ha_test:5, ha_test:33}
\end{requirement}

\begin{requirement}{R-0064}{A generation is never issued twice, and never goes backwards}
  Each grant of entitlement shall carry a generation greater than every
  generation the venue has used. Where every instance that could have issued one
  has restarted, the new generation shall still exceed every generation
  previously used, which requires that the number outlive the processes rather
  than any one of them.
  \rationale{The generation is what lets a receiver refuse a superseded sender. A
  generation held only in memory restarts at the beginning, so a pair that
  restarts together agrees a generation the venue has used before, and a message
  left over from the earlier one then compares as current.

  This is the requirement most exposed to a restart of everything at once. Where
  each instance keeps its own record, a group whose members all restart can still
  exceed every generation used, because each instance reads its own record and
  the highest wins. Where the record is kept in one place, the survival of that
  place is the whole of the guarantee. This chapter does not settle which of the
  two the venue uses. The requirement holds either way, and one of them has to be
  chosen deliberately.}
  \uncovered{no scenario restarts every instance of a group together}
\end{requirement}

\begin{gap}{}
  A generation can go backwards. An arbiter that has restarted, and has not been
  told the highest generation by the other arbiter because it has restarted too,
  can grant a new leader a generation below one already led in. Two instances
  still never act at once, but receivers that saw the higher generation refuse
  what the new leader sends. The new leader learns of the higher generation when
  its peer refuses it, and asks again above it. Where the peer is also down, only
  a receiver could tell it, and receivers do not yet do so. See
  \texttt{docs/availability/tla/findings.md}, section 11.5.
\end{gap}

\begin{requirement}{R-0065}{An instance whose entitlement has passed has what it sends refused}
  A receiver shall refuse a message carrying a generation older than the newest
  it has accepted for that group.
  \rationale{An instance whose lease has run out stops acting by itself
  (\reqref{R-0146}), before anything it sends could be refused. Refusing an older
  generation is the second defence, for anything that reaches a receiver all the
  same: a message delayed in transit, or a leader granted a generation below one a
  receiver has already seen. It needs no special hardware and no cooperation from
  the instance being refused.}
  \uncovered{a superseded instance stops before it sends, so no scenario makes one send anything for a receiver to refuse}
\end{requirement}
